% Options for packages loaded elsewhere
\PassOptionsToPackage{unicode}{hyperref}
\PassOptionsToPackage{hyphens}{url}
\documentclass[
  12pt,
  a4paper,
]{article}
\usepackage{xcolor}
\usepackage[margin=1in]{geometry}
\usepackage{amsmath,amssymb}
\setcounter{secnumdepth}{-\maxdimen} % remove section numbering
\usepackage{iftex}
\ifPDFTeX
  \usepackage[T1]{fontenc}
  \usepackage[utf8]{inputenc}
  \usepackage{textcomp} % provide euro and other symbols
\else % if luatex or xetex
  \usepackage{unicode-math} % this also loads fontspec
  \defaultfontfeatures{Scale=MatchLowercase}
  \defaultfontfeatures[\rmfamily]{Ligatures=TeX,Scale=1}
\fi
\usepackage{lmodern}
\ifPDFTeX\else
  % xetex/luatex font selection
\fi
% Use upquote if available, for straight quotes in verbatim environments
\IfFileExists{upquote.sty}{\usepackage{upquote}}{}
\IfFileExists{microtype.sty}{% use microtype if available
  \usepackage[]{microtype}
  \UseMicrotypeSet[protrusion]{basicmath} % disable protrusion for tt fonts
}{}
\makeatletter
\@ifundefined{KOMAClassName}{% if non-KOMA class
  \IfFileExists{parskip.sty}{%
    \usepackage{parskip}
  }{% else
    \setlength{\parindent}{0pt}
    \setlength{\parskip}{6pt plus 2pt minus 1pt}}
}{% if KOMA class
  \KOMAoptions{parskip=half}}
\makeatother
% definitions for citeproc citations
\NewDocumentCommand\citeproctext{}{}
\NewDocumentCommand\citeproc{mm}{%
  \begingroup\def\citeproctext{#2}\cite{#1}\endgroup}
\makeatletter
 % allow citations to break across lines
 \let\@cite@ofmt\@firstofone
 % avoid brackets around text for \cite:
 \def\@biblabel#1{}
 \def\@cite#1#2{{#1\if@tempswa , #2\fi}}
\makeatother
\newlength{\cslhangindent}
\setlength{\cslhangindent}{1.5em}
\newlength{\csllabelwidth}
\setlength{\csllabelwidth}{3em}
\newenvironment{CSLReferences}[2] % #1 hanging-indent, #2 entry-spacing
 {\begin{list}{}{%
  \setlength{\itemindent}{0pt}
  \setlength{\leftmargin}{0pt}
  \setlength{\parsep}{0pt}
  % turn on hanging indent if param 1 is 1
  \ifodd #1
   \setlength{\leftmargin}{\cslhangindent}
   \setlength{\itemindent}{-1\cslhangindent}
  \fi
  % set entry spacing
  \setlength{\itemsep}{#2\baselineskip}}}
 {\end{list}}
\usepackage{calc}
\newcommand{\CSLBlock}[1]{\hfill\break\parbox[t]{\linewidth}{\strut\ignorespaces#1\strut}}
\newcommand{\CSLLeftMargin}[1]{\parbox[t]{\csllabelwidth}{\strut#1\strut}}
\newcommand{\CSLRightInline}[1]{\parbox[t]{\linewidth - \csllabelwidth}{\strut#1\strut}}
\newcommand{\CSLIndent}[1]{\hspace{\cslhangindent}#1}
\setlength{\emergencystretch}{3em} % prevent overfull lines
\providecommand{\tightlist}{%
  \setlength{\itemsep}{0pt}\setlength{\parskip}{0pt}}
\usepackage{bookmark}
\IfFileExists{xurl.sty}{\usepackage{xurl}}{} % add URL line breaks if available
\urlstyle{same}
\hypersetup{
  hidelinks,
  pdfcreator={LaTeX via pandoc}}

\author{}
\date{}

\begin{document}

\section{Introduction to Digital Twins for Software
Engineers}\label{introduction-to-digital-twins-for-software-engineers}

\subsection{Learning objectives}\label{learning-objectives}

By the end of this chapter you will be able to:

\begin{enumerate}
\def\labelenumi{\arabic{enumi}.}
\tightlist
\item
  \textbf{Classify} a proposed system as a digital model, a digital
  shadow, or a digital twin, by tracing the direction of its data flows.
\item
  \textbf{Identify}, in a running system, the three pieces of state
  every digital twin holds: what it sensed, what it believes, and what
  it commanded.
\item
  \textbf{Predict} how a twin misbehaves once its model stops matching
  the physical object.
\item
  \textbf{Derive} a check that detects that divergence from sensor
  readings alone.
\end{enumerate}

\subsection{Scope and prerequisites}\label{scope-and-prerequisites}

One worked example runs from start to finish: a potted plant, a
soil-moisture sensor, and a pump that waters it. Through it, the chapter
covers the model/shadow/twin distinction, the anatomy of a twin, and the
way twins fail as their models age. It does not cover co-simulation,
modelling formalisms, or any specific commercial platform, and it is
explanation rather than a build guide --- to type code and watch a twin
run, work through \texttt{digital-twin-plant-moisture-tutorial.md}
alongside it.

You are assumed to write software professionally and to be comfortable
with state, scheduled jobs, and testing. You are not assumed to know
control theory, or ever to have wired up a sensor.

\subsection{1. Why an engineer should
care}\label{why-an-engineer-should-care}

You already know how to build a service that reads a sensor and draws a
chart. A digital twin is the harder thing that chart is often mistaken
for.

The distinction has teeth because the twin \emph{acts}. A dashboard that
is wrong produces a confused operator. A twin that is wrong waters the
plant at 3 a.m. until the pot overflows. Bidirectional interaction
between the physical object and its virtual counterpart is the centre of
the concept, not an optional extra {[}1{]}. The discipline is also
unsettled: practitioners report that the techniques for building and
maintaining twins are still not well understood {[}2{]}. You are not
arriving late to a solved problem.

Four terms, defined once and used consistently from here on. The
\textbf{physical twin} is the real object: this pot, this plant. The
\textbf{digital twin} is the software that mirrors it. A \textbf{sensor}
reports the physical twin's state to the software; an \textbf{actuator}
changes it on the software's command. The pump is an actuator.

\subsection{2. Three systems, one plant}\label{three-systems-one-plant}

Kritzinger and colleagues separate three arrangements by how much of the
data flow is automatic {[}3{]}. Take them in order --- the sequence is
the lesson.

\textbf{A digital model.} You keep a spreadsheet holding a drying-rate
estimate for your plant. Once a week you poke the soil, type in a
number, and the sheet predicts when to water next. No data moves on its
own in either direction; a person carries every reading in and every
decision out.

\textbf{A digital shadow.} Now you push a moisture probe into the pot.
It posts a reading every fifteen minutes to a service that stores it and
plots it. Data now flows automatically from pot to software --- but only
that way. When the chart dips, \emph{you} fill the watering can. Most
systems sold as digital twins stop here.

\textbf{A digital twin.} The same service now opens a valve when it
decides the soil is too dry. Data flows automatically in both
directions, and the loop closes without you. That closure is the whole
difference, and notice what it costs: a wrong shadow annoys you, while a
wrong twin has a pump.

So the first question to ask about any system claiming to be a twin is
not what it models. It is: \textbf{what does it change, and who approved
that change?}

\subsection{3. Anatomy: what the twin actually
holds}\label{anatomy-what-the-twin-actually-holds}

Let \(m(t)\) be the soil moisture at hour \(t\), as a percentage of the
pot's capacity, and let \(r\) be the drying rate in percentage points
per hour. The twin's model is one line:

\[m(t + 1) = m(t) - r\]

Suppose the probe reads \(m(0) = 60\), the twin's stored estimate is
\(r = 2.0\), and the plant wilts below 30. The twin projects forward:

\[t_{\text{dry}} = \frac{60 - 30}{2.0} = 15 \text{ hours}\]

Three distinct pieces of state are now in play, and confusing any two of
them is the classic beginner's bug:

\begin{itemize}
\tightlist
\item
  \textbf{Sensed state} --- 60. What the probe said. Noisy, and already
  stale by the time it arrives.
\item
  \textbf{Modelled state} --- the estimate \(r = 2.0\) and the
  projection to 15 hours. What the twin \emph{believes}. This exists
  nowhere in the physical world.
\item
  \textbf{Commanded state} --- ``open the valve for 8 seconds at hour
  15''. What the twin decided to do about the gap.
\end{itemize}

A dashboard has only the first. A twin has all three, and must reconcile
them every cycle, because the physical twin does not know what the
software believes.

\subsection{4. How this breaks}\label{how-this-breaks}

The plant grows. Summer arrives. The real drying rate climbs to
\(r = 3.0\), and the soil crosses 30 at hour 10. The twin, still holding
\(r = 2.0\), waters at hour 15. The plant sat wilting for five hours
while the software reported that all was well --- and the twin's logs
show no error at all. Nothing crashed. The model aged out from under the
plant, and nothing in the system was built to notice.

This gap between a twin's model and its physical twin is why serious
work in this area treats \textbf{online model updating} --- refitting a
model from data collected while the system runs --- as a core
requirement rather than a refinement {[}4{]}.

Work the detection out one step at a time. The twin predicted
\(m(1) = 58\); the probe reported 57. One point of disagreement is
sensor noise, not evidence. By hour six, though, the twin predicts 48
and the probe says 42. \textbf{The residual --- prediction minus
measurement --- is growing in one direction.} That signature is what you
watch for: noise cancels out over time, while a one-sided bias means the
model itself is wrong.

The remaining step is yours (Exercise 2): decide what the twin should do
about it.

\subsection{5. Exercises}\label{exercises}

\textbf{Exercise 1 (objective 1).} Your neighbour's system texts them
``water me'', and they do. Classify it, then state the smallest change
that would make it a digital twin. \emph{Hint: count the automatic flows
and their directions.}

\textbf{Exercise 2 (objectives 3, 4).} Continuing §4: the twin has six
hours of one-sided residuals. Write down a rule that re-estimates \(r\)
from the last \(N\) readings. Then the harder half --- what happens if
\(N\) is small and a spilled cup of tea enters the window? \emph{Hint:
the rule must not trust one reading more than it trusts the model.}

\textbf{Exercise 3 (objective 2).} The valve sticks open. Describe how
sensed, modelled and commanded state now disagree, and design the check
that catches it. \emph{Hint: after watering, the twin makes a prediction
it can verify.}

\subsection{6. Summary and where to go
next}\label{summary-and-where-to-go-next}

A digital twin is not a model, and not a dashboard. It is a model wired
to a physical object in both directions, so that it senses, believes,
and acts (§2, §3). Its characteristic failure is not a crash but silent
divergence: the model ages, the object changes, and the twin keeps
acting confidently on a stale belief (§4).

For the practices around building one, Talasila and colleagues work
through what realising a twin from real assets and models demands
{[}5{]}, and Gil and colleagues compare open-source frameworks worth
building on rather than reinventing {[}6{]}. The companion tutorial
builds this plant twin in about seventy lines of Python.

\subsection{7. References}\label{references}

\protect\phantomsection\label{refs}
\begin{CSLReferences}{0}{0}
\bibitem[\citeproctext]{ref-committee_on_foundational_research_gaps_and_future_directions_for_digital_twins_foundational_2024}
\CSLLeftMargin{{[}1{]} }%
\CSLRightInline{\emph{Foundational {Research} {Gaps} and {Future}
{Directions} for {Digital} {Twins}}. Washington, D.C.: National
Academies Press, 2024. doi:
\href{https://doi.org/10.17226/26894}{10.17226/26894}.}

\bibitem[\citeproctext]{ref-cleophas_community-sourced_2022}
\CSLLeftMargin{{[}2{]} }%
\CSLRightInline{L. Cleophas \emph{et al.}, {``A community-sourced view
on engineering digital twins: A report from the {EDT}.{Community},''} in
\emph{Proceedings of the 25th {International} {Conference} on {Model}
{Driven} {Engineering} {Languages} and {Systems}: {Companion}
{Proceedings}}, in {MODELS} '22. New York, NY, USA: Association for
Computing Machinery, Nov. 2022, pp. 481--485. doi:
\href{https://doi.org/10.1145/3550356.3561549}{10.1145/3550356.3561549}.}

\bibitem[\citeproctext]{ref-kritzinger_digital_2018}
\CSLLeftMargin{{[}3{]} }%
\CSLRightInline{W. Kritzinger, M. Karner, G. Traar, J. Henjes, and W.
Sihn, {``Digital {Twin} in manufacturing: {A} categorical literature
review and classification,''} \emph{Ifac-PapersOnline}, vol. 51, no. 11,
pp. 1016--1022, 2018, Accessed: Dec. 11, 2024. {[}Online{]}. Available:
\url{https://www.sciencedirect.com/science/article/pii/S2405896318316021}}

\bibitem[\citeproctext]{ref-thelen_comprehensive_2022}
\CSLLeftMargin{{[}4{]} }%
\CSLRightInline{A. Thelen \emph{et al.}, {``A comprehensive review of
digital twin~---~part 1: Modeling and twinning enabling technologies,''}
\emph{Structural and Multidisciplinary Optimization}, vol. 65, no. 12,
p. 354, Nov. 2022, doi:
\href{https://doi.org/10.1007/s00158-022-03425-4}{10.1007/s00158-022-03425-4}.}

\bibitem[\citeproctext]{ref-talasila_realising_2024}
\CSLLeftMargin{{[}5{]} }%
\CSLRightInline{P. Talasila, P. H. Mikkelsen, S. Gil, and P. G. Larsen,
{``Realising {Digital} {Twins},''} in \emph{The {Engineering} of
{Digital} {Twins}}, J. Fitzgerald, C. Gomes, and P. G. Larsen, Eds.,
Cham: Springer International Publishing, 2024, pp. 225--256. doi:
\href{https://doi.org/10.1007/978-3-031-66719-0_11}{10.1007/978-3-031-66719-0\_11}.}

\bibitem[\citeproctext]{ref-gil_survey_2024}
\CSLLeftMargin{{[}6{]} }%
\CSLRightInline{S. Gil, P. H. Mikkelsen, C. Gomes, and P. G. Larsen,
{``Survey on open‐source digital twin frameworks--{A} case study
approach,''} \emph{Software: Practice and Experience}, vol. 54, no. 6,
pp. 929--960, Jun. 2024, doi:
\href{https://doi.org/10.1002/spe.3305}{10.1002/spe.3305}.}

\end{CSLReferences}

\end{document}
